% !TEX root =  master.te
\chapter{Grundlagen}

\section{Audit}

Das Audit stellt eine unabhänige Prüfungsfunktion innerhalb einer Organisation dar. Die zentrale Aufgabe besteht darin die angemessene Wirksamkeit von Governance- und internen Kontrollsystemen zu bewerten und hierzu objektive Prüfungs und Beratungsleistungen bereitzustellen. (vgl. \cite[S. 263 - 264]{cantu2025making}). Dafür müssen Dokumente gesichtet, interpretiert und geprüft werden (vgl. \cite[S.556 -557]{ROUSSY2013550}). Im Rahmen dieser Arbeit liegt der Fokus auf dem Internal Audit bei \ac{FCo}, bei der Dokumente aus verschiedenen Konzernfunktionen ausgewertet werden.



\section{Stichprobenanalyse im Audit}
Für das Audit ist es in der Regel nicht möglich, alle Dokumente einer Grundgesamtheit zu prüfen. Daher wird auf die Stichprobenanalyse zurückgegriffen. Sie ist eine Prüfungsmethode im Audit-Bereich, bei der eine möglichst hohe Prüfungseffizienz erreicht werden soll. Gewählt wird demnach diejenige Handlungsalternative, die mit dem geringsten Mitteleinsatz zu einem hinreichend sicheren Prüfungsurteil führt. Aus einer Grundgesamtheit, also einer abgegrenzten Menge gleichartiger Elemente, die beispielsweise zu einem Konto gehören, werden hierzu eine Pflichtprobe und eine zufällige Stichprobe entnommen.
Die Pflichtprobe umfasst Elemente, die zuvor definierte Kriterien erfüllen, beispielsweise Rechnungen, die über einem bestimmten Betrag liegen (vgl. \cite{freudenberg2026auditsample}). Diese Elemente müssen als Teil der Stichprobe geprüft werden. Weitergehend wird aus den restlichen Elementen der Grundgesamtheit eine zufällige Stichprobe entnommen, die zusätzlich zur Pflichtprobe geprüft wird. Die verbleibenden Elemente der Grundgesamtheit werden vorerst nicht geprüft und bedingen damit eine gewisse Unschärfe der Prüfungsaussagen, die als Stichprobenrisiko bezeichnet wird (vgl. \cite[S.\,11--12]{baumeister2019sample}).

\section{Künstliche Intelligenz}

\subsection{Large Language Models}

\subsection{Prompt Engineering}

Das Prompt Engineering beschreibt den wichitgsten Part der Nutzung eines KI-Systems. Erst durch den Prompt weiß die Maschine was der User erreichen möchte. Ein Prompt kann in verschiedenen klar definierte Teile unterteilt sein.



\section{Einsatz von Künstlicher Intelligenz im Audit}

\subsection{Potenziale in der Stichprobenanalyse}

\subsection{Chancen und Risiken}

\subsection{KI-gestütze Audit-Lösungen}

DataSnipper